\subsection{Relations: array, collection, one}
\label{sec:go-relations}

A dto's \texttt{one}/\texttt{collection} field just needs a Go type to
hold data. An entity's relation field needs that \emph{and} a real,
migratable gorm association --- and the two shapes are incompatible in a
way that isn't obvious up front. \texttt{array}/\texttt{collection}/
\texttt{one} normally render as \texttt{emigo.Array[T]}/
\texttt{Collection[T]}/\texttt{One[T]} --- PATCH-payload wrapper types
carrying an \texttt{Operation} (\texttt{replace}/\texttt{append}/
\texttt{select}) --- the right shape for an API, but not one gorm can
migrate (verified directly against \texttt{gorm.io/gorm/schema.Parse}).

So an entity keeps the DTO field as-is (still the only public JSON/CLI
surface) tagged \texttt{gorm:"-"}, and generates a hidden sibling field
that \emph{is} a real gorm association:

\begin{lstlisting}
Items    emigo.Array[Entity1EntityItems] `gorm:"-" json:"items"`
Owner    emigo.One[Entity2Entity]        `gorm:"-" json:"owner"`

ItemsRow []*Entity1EntityItems `gorm:"foreignKey:LinkerId;references:Id;
                                 constraint:OnDelete:CASCADE" json:"-"`
OwnerId  int64          `gorm:"index" json:"-"`
OwnerRow *Entity2Entity `gorm:"foreignKey:OwnerId;references:Id" json:"-"`
\end{lstlisting}
The \texttt{\{field\}Row}/\texttt{\{field\}Id} siblings are purely a
persistence concern (\texttt{json:"-"}); whatever writes to the database
(\S\ref{sec:go-actions}) keeps the DTO field and its sibling in sync.

\begin{sidebar}{The three relation kinds}
\textbf{array} --- has-many, owned composition. Nested \texttt{fields:}
become a child struct with its own \texttt{Id}/\texttt{UniqueId} plus a
\texttt{LinkerId} FK back to the parent.

\textbf{collection} --- many-to-many, backed by a join table named
\texttt{\{entity\}\_\{field\}} (e.g.\ \texttt{entity1\_items3}).

\textbf{one} --- belongs-to: a foreign-key column on \emph{this} entity,
named \texttt{\{field\}Id}, referencing the target's \texttt{Id}.
\end{sidebar}

\textbf{Why \texttt{Id} and not \texttt{UniqueId}}: every
\texttt{foreignKey}/\allowbreak\texttt{references} tag joins on the small integer
\texttt{Id}, keeping indexes smaller and joins faster, even though
\texttt{UniqueId} is the only identity an API caller ever actually has.
Since gorm's \texttt{Save()} decides insert-vs-update purely from whether
\texttt{Id} is already populated, and a caller only ever knows
\texttt{UniqueId}, the reconcile helpers (\S\ref{sec:go-actions}) resolve
an incoming item's real \texttt{Id} by matching \texttt{UniqueId} first.

\subsubsection{Relations nested inside object containers}
A relation doesn't have to sit directly on the entity --- it can live
inside an \texttt{object} (or \texttt{object?}) field, any depth deep.
Gorm's own embedding correctly walks an arbitrarily deep \texttt{embedded}
chain and still recognizes the association, keyed by its bare field name
rather than the full nested path (verified against
\texttt{gorm.io/gorm/schema.Parse} and a real database run) --- provided
the nested child struct is named with the same accumulating prefix the
struct generator uses for the container itself, and
\texttt{\{Entity\}CreateFn}/\texttt{UpdateFn} recurse into
\texttt{object}/\texttt{object?} containers to find relations at any
depth. \texttt{object?} specifically needs its own fix: it renders as
\texttt{emigo.Nullable[T]}, and gorm's embedding only ever sees the
wrapper's own unexported fields, never \texttt{T}'s --- so it gets the
same \texttt{\{field\}Row *T} hidden-sibling treatment a relation does,
even for a field with no relation in it at all.

\begin{pullquote}
Known limitations: a relation declared on an array item itself (rather
than on the entity or one of its object containers) isn't picked up ---
each array item is still persisted as one flat unit via \texttt{Save()}.
Cross-module targets (\texttt{module:} on a relation field) aren't
accounted for in the FK/table naming yet.
\end{pullquote}
